\documentclass[11pt]{article}

\usepackage[T1]{fontenc}
\usepackage[utf8]{inputenc}
\usepackage{lmodern}
\usepackage[margin=1in]{geometry}
\usepackage{microtype}
\usepackage{amsmath,amssymb,amsthm,mathtools}
\usepackage{xcolor}
\usepackage[numbers,sort&compress]{natbib}
\usepackage{hyperref}

\hypersetup{
  colorlinks=true,
  linkcolor=blue!55!black,
  citecolor=green!40!black,
  urlcolor=blue!65!black,
  pdftitle={Fault-Tolerant Homomorphic Measurement for Lifted Product Codes},
  pdfauthor={Research Progress}
}

\newtheorem{definition}{Definition}[section]
\newcommand{\F}{\mathbb{F}}
\newcommand{\B}{\mathsf{B}}
\newcommand{\coker}{\operatorname{coker}}
\newcommand{\im}{\operatorname{im}}
\newcommand{\Supp}{\operatorname{Supp}}

\title{Fault-Tolerant Homomorphic Measurement for Lifted Product Codes}
\author{Research Progress}
\date{October 8, 2026}

\begin{document}

\maketitle

\begin{abstract}
For CSS codes, homomorphic measurement provides a general framework for
measuring logical Pauli operators by coupling the data block to an ancillary
code through a suitable code homomorphism. However, only a limited number of
code families currently have both fault-tolerant ancilla constructions and
explicit homomorphisms that support addressable logical measurements. In this
work, we investigate fault-tolerant constructions for addressable homomorphic
measurement in lifted product codes. Our goal is to identify ancillary lifted
product codes, compatible homomorphisms, and fault-tolerant preparation and
readout procedures that implement selected logical measurements while
preserving the remaining logical information.
\end{abstract}

\section{Background}

\subsection{CSS codes as three-term chain complexes}

A CSS code is represented by a three-term chain complex over $\F_2$,
\begin{equation}
  Q:\qquad
  Q_2 \xrightarrow{\,\partial_2^Q\,}
  Q_1 \xrightarrow{\,\partial_1^Q\,} Q_0,
  \qquad
  \partial_1^Q\partial_2^Q=0.
  \label{eq:css-complex}
\end{equation}
The bases of $Q_2$, $Q_1$, and $Q_0$ index the $Z$ checks, qubits, and $X$
checks, respectively. Thus
\begin{equation}
  H_X=\partial_1^Q,
  \qquad
  H_Z^T=\partial_2^Q,
  \qquad
  H_XH_Z^T=0.
  \label{eq:css-checks}
\end{equation}
Logical $Z$ operators are homology classes and logical $X$ operators are
cohomology classes,
\begin{equation}
  H_1(Q)=\frac{\ker\partial_1^Q}{\im\partial_2^Q},
  \qquad
  H^1(Q)=
  \frac{\ker(\partial_2^Q)^T}{\im(\partial_1^Q)^T}.
  \label{eq:logical-homology}
\end{equation}
\subsection{Lifted-product codes}

\subsubsection{The ring-valued product complex}

Let
\begin{equation}
  R=\F_2[x]/(x^\ell+1).
  \label{eq:ring}
\end{equation}
Over $\F_2$, $x^\ell+1=x^\ell-1$, so multiplication by $x$ is the cyclic
shift of a length-$\ell$ binary fibre. For $f\in R$, write
$\overline f=f(x^{-1})=f(x^{\ell-1})$, and for a matrix $C$ over $R$ define
$C^*=\overline C^{\,T}$. Consider two one-term complexes of free $R$-modules,
\begin{equation}
  A_\bullet:\ R^{n_A}\xrightarrow{\,A\,}R^{m_A},
  \qquad
  B_\bullet:\ R^{n_B}\xrightarrow{\,B\,}R^{m_B},
  \label{eq:base-complexes}
\end{equation}
where $A\in R^{m_A\times n_A}$ and
$B\in R^{m_B\times n_B}$. The two-dimensional lifted product
$\operatorname{LP}_\ell(A,B)$ is the total complex $M=A_\bullet\otimes_R
B_\bullet$,
\begin{equation}
  M_2 \xrightarrow{\,\partial_2^M\,}
  M_1 \xrightarrow{\,\partial_1^M\,} M_0,
  \label{eq:lp-complex}
\end{equation}
with
\begin{align}
  M_2 &= R^{n_A}\otimes_R R^{n_B}, \notag\\
  M_1 &=
  \bigl(R^{n_A}\otimes_R R^{m_B}\bigr)
  \oplus
  \bigl(R^{m_A}\otimes_R R^{n_B}\bigr),
  \label{eq:lp-modules}\\
  M_0 &= R^{m_A}\otimes_R R^{m_B}, \notag
\end{align}
and boundary maps
\begin{equation}
  (\partial_2^M)^T=
  \bigl(I_{n_A}\otimes_R B^*,\ A^*\otimes_R I_{n_B}\bigr),
  \qquad
  \partial_1^M=
  \bigl(A\otimes_R I_{m_B},\ I_{m_A}\otimes_R B\bigr).
  \label{eq:lp-boundaries}
\end{equation}
These are the conventions of Definition III.1 in
Ref.~\cite{ZhengEtAl2026}. In particular,
$\partial_1^M\partial_2^M=0$.

\subsubsection{Binarization and the CSS code}

The monomial basis $1,x,\ldots,x^{\ell-1}$ defines an $\F_2$-linear
isomorphism
\begin{equation}
  \B:R^s\longrightarrow\F_2^{s\ell}.
  \label{eq:binarization}
\end{equation}
It also gives the regular binary representation of an $R$-linear matrix and
satisfies $\B(Cu)=\B(C)\B(u)$. Applying $\B$ to
Eq.~\eqref{eq:lp-complex} gives a binary chain complex $Q=\B(M)$ and hence a
CSS code with
\begin{equation}
  H_X=\B(\partial_1^M),
  \qquad
  H_Z^T=\B(\partial_2^M).
  \label{eq:lp-css-checks}
\end{equation}
The physical-qubit count is therefore
\begin{equation}
  n=\ell\bigl(n_A m_B+m_A n_B\bigr).
  \label{eq:lp-length}
\end{equation}
The map $\B$ is a bijective chain map, so it identifies the ring-valued
homology of $M$ with the binary logical-$Z$ space of $Q$. When $\ell$ is odd,
$R$ is semisimple and the K\"unneth decomposition takes the schematic form
\begin{equation}
  H_1(M)\simeq
  \bigl(\ker_R A\otimes_R\coker_R B\bigr)
  \oplus
  \bigl(\coker_R A\otimes_R\ker_R B\bigr).
  \label{eq:kunneth}
\end{equation}
This decomposition explains the two physical sectors in $M_1$ and is the
starting point for constructing explicit logical bases.

\subsubsection{Canonical logical fibres}

Zheng \emph{et al.} identify a canonical sector when compatible information
sets can be chosen for the relevant kernel and cokernel modules. Put
$r_A=n_A-m_A$ and $r_B^*=m_B-n_B$. Its logical qubits are indexed by
\begin{equation}
  (i,j,m)\in [r_A]\times[r_B^*]\times\mathbb Z_\ell,
  \qquad
  k_{\mathrm c}=r_A r_B^*\ell.
  \label{eq:canonical-index}
\end{equation}
For suitable systematic representatives
$u_{A_1}^{(i)},e_{A_1}^{(i)},e_{B_0}^{(j)},v_{B_0}^{(j)}$, a conjugate logical
pair has ring-valued supports
\begin{align}
  \Supp\overline Z_{i,j,m}
  &=x^m u_{A_1}^{(i)}\otimes_R e_{B_0}^{(j)},
  \label{eq:canonical-z}\\
  \Supp\overline X_{i,j,m}
  &=x^m e_{A_1}^{(i)}\otimes_R v_{B_0}^{(j)}.
  \label{eq:canonical-x}
\end{align}
Their binary supports are obtained with $\B$. The $Z$ and $X$ representatives
intersect on exactly one physical qubit precisely when all three indices agree.
Multiplication by $x^t$ shifts the third index,
$m\mapsto m+t\pmod\ell$. Thus an HGP-like row--column grid is replaced by a
grid of length-$\ell$ logical fibres. The canonical sector need not exhaust
the full logical space; any residual logical sector must be tracked separately
when claiming addressability or nondestructive action on spectator qubits.

\subsection{Homomorphic CNOT and homomorphic measurement}

\subsubsection{Chain maps as physical CNOT patterns}

Let $Q$ be a data code and $Q'$ an ancilla code, both written as in
Eq.~\eqref{eq:css-complex}. A chain map $\gamma:Q'\to Q$ is a triple
\begin{equation}
  \gamma_i:Q_i'\longrightarrow Q_i,
  \qquad i=0,1,2,
  \label{eq:chain-map-components}
\end{equation}
satisfying
\begin{equation}
  \partial_i^Q\gamma_i
  =\gamma_{i-1}\partial_i^{Q'},
  \qquad i=1,2.
  \label{eq:chain-map}
\end{equation}
The map $\gamma_1$ specifies the physical CNOTs: data qubit $i$ controls
ancilla qubit $j$ whenever $\gamma_1[i,j]=1$. Equation~\eqref{eq:chain-map}
ensures stabilizer preservation. Coordinate inclusions are transversal, while
bounded-degree maps have bounded CNOT depth.

The logical action is most transparent on (co)homology. The chain map induces
\begin{equation}
  \gamma_*:H_1(Q')\longrightarrow H_1(Q),
  \qquad
  \gamma^*:H^1(Q)\longrightarrow H^1(Q').
  \label{eq:induced-maps}
\end{equation}
For a data-control, ancilla-target CNOT, an ancilla logical-$Z$ class $[z']$
and a data logical-$X$ class $[x]$ transform schematically as
\begin{align}
  \overline Z_{Q'}([z'])
  &\longmapsto
  \overline Z_Q(\gamma_*[z'])\,\overline Z_{Q'}([z']),
  \label{eq:z-action}\\
  \overline X_Q([x])
  &\longmapsto
  \overline X_Q([x])\,\overline X_{Q'}(\gamma^*[x]).
  \label{eq:x-action}
\end{align}
The two induced maps therefore have to be checked together: specifying the
desired image of a $Z$ logical does not by itself establish the full logical
Clifford action.

\subsubsection{Generalized Steane measurement}

The homomorphic CNOT becomes a nondestructive logical measurement through a
generalized Steane circuit. For a $Z$-type measurement, prepare the relevant
logical qubits of $Q'$ in the logical $Z$ basis, apply the data-control
homomorphic CNOT, and measure $Q'$ in the physical $Z$ basis. Each decoded
ancilla logical outcome measures the data logical product represented by the
corresponding class in $\im\gamma_*$. For an $X$-type measurement, prepare
the ancilla in the logical $X$ basis, reverse the physical CNOT direction, and
measure the ancilla in the physical $X$ basis. Multiple ancilla logical
qubits give multiple commuting Pauli-product measurements in parallel.

A chain map proves that the ideal CNOT pattern preserves the codespace, but it
is not by itself a fault-tolerance proof. The construction in
Ref.~\cite{XuEtAl2024} additionally uses an ancilla whose distance is at least
the data distance, fault-tolerant logical-state preparation, a constant-depth
or transversal coupling, repeated syndrome extraction, transversal readout,
and classical decoding or Pauli-frame processing. These are separate
obligations in an LP extension.

\subsection{HGP homomorphic CNOTs}

\subsubsection{HGP codes inside the LP convention}

Let $H_a\in\F_2^{r_a\times n_a}$ be full-row-rank classical check matrices,
with $k_a=n_a-r_a$ for $a=1,2$. Xu \emph{et al.} use $H_1^T$ as the vertical
chain map and $H_2$ as the horizontal chain map. In the LP notation above,
this is the specialization
\begin{equation}
  \ell=1,
  \qquad
  A=H_1^T,
  \qquad
  B=H_2.
  \label{eq:hgp-lp-translation}
\end{equation}
After swapping the two direct summands of $M_1$, the resulting HGP check
matrices are
\begin{align}
  H_X&=\bigl(I_{n_1}\otimes H_2,\ H_1^T\otimes I_{r_2}\bigr),
  \label{eq:hgp-hx}\\
  H_Z&=\bigl(H_1\otimes I_{n_2},\ I_{r_1}\otimes H_2^T\bigr).
  \label{eq:hgp-hz}
\end{align}
The code has parameters
\begin{equation}
  [[n_1n_2+r_1r_2,\ k_1k_2,\ \min(d_1,d_2)]],
  \label{eq:hgp-parameters}
\end{equation}
under the full-rank assumptions used in Ref.~\cite{XuEtAl2024}.

Choose systematic kernel generators $b_i^{(1)}\in\ker H_1$ and
$b_j^{(2)}\in\ker H_2$, together with unit vectors on the corresponding
information coordinates. A canonical logical grid is
\begin{equation}
  \overline Q=\{\overline Q_{i,j}\}_{i\in[k_1],j\in[k_2]},
  \label{eq:hgp-grid}
\end{equation}
with representatives, in the bit--bit qubit sector,
\begin{equation}
  \overline X_{i,j}=
  \begin{pmatrix}b_i^{(1)}\otimes e_j^{(n_2)}\\0\end{pmatrix},
  \qquad
  \overline Z_{i,j}=
  \begin{pmatrix}e_i^{(n_1)}\otimes b_j^{(2)}\\0\end{pmatrix}.
  \label{eq:hgp-logicals}
\end{equation}
The $X$ representative occupies a logical column, the $Z$ representative a
logical row, and conjugate representatives meet at one information-coordinate
qubit.

\subsubsection{Puncturing, augmenting, and the induced CNOT}

Consider a collection of disjoint horizontal hyperedges
$E_h=\{e_\alpha\}$ with $e_\alpha\subseteq[k_2]$. Let
\begin{equation}
  S=[k_2]\setminus\bigcup_{e\in E_h}e
  \label{eq:puncture-set}
\end{equation}
be the unused horizontal information coordinates. Starting from $H_2$, add
repetition-code checks $H_{\mathrm{rep}}^{E_h}$ that identify the logical bits
inside each $e_\alpha$, and then puncture the coordinates in $S$:
\begin{equation}
  H_2'=\bigl(H_2^{+H_{\mathrm{rep}}^{E_h}}\bigr)^S,
  \qquad
  Q'=\operatorname{HGP}(H_1,H_2').
  \label{eq:hgp-ancilla}
\end{equation}
Augmentation expurgates unwanted codewords, while puncturing information
coordinates removes unused logical columns. The coordinate inclusion from
$H_2'$ to $H_2$, tensored with the identity on the vertical factor, induces a
chain map $\gamma:Q'\to Q$. Its qubit component identifies the shared HGP
coordinates, so the physical homomorphic CNOT is transversal on those
coordinates.

For each row $i$ and hyperedge $e_\alpha$, one logical qubit of $Q'$ couples to
all data logical qubits $\{\overline Q_{i,j}:j\in e_\alpha\}$. Up to
stabilizers, the data-control, ancilla-target CNOT acts as
\begin{align}
  \overline Z'_{i,\alpha}
  &\longmapsto
  \overline Z'_{i,\alpha}
  \prod_{j\in e_\alpha}\overline Z_{i,j},
  \label{eq:hgp-logical-z-action}\\
  \overline X_{i,j}
  &\longmapsto
  \overline X_{i,j}\overline X'_{i,\alpha},
  \qquad j\in e_\alpha,
  \label{eq:hgp-logical-x-action}
\end{align}
while the remaining logical operators are unchanged. Preparing $Q'$ in the
logical $Z$ basis and measuring it after this CNOT therefore measures
\begin{equation}
  \left\{
  \prod_{j\in e}\overline Z_{i,j}
  : i\in[k_1],\ e\in E_h
  \right\}
  \label{eq:horizontal-ppms}
\end{equation}
in parallel. This is the horizontal-PPM gadget of Algorithm 1 and Theorem 9
in Ref.~\cite{XuEtAl2024}.

The fault-tolerance argument uses two special properties of this construction.
First, the punctured coordinates are information bits, and the added checks
only expurgate codewords; consequently the modified horizontal classical
distance does not decrease. Second, $Q'$ is prepared by transversal physical
$|0\rangle$ preparation followed by $d=\min(d_1,d_2)$ rounds of stabilizer
measurement, coupled transversally to $Q$, and read out transversally. Thus the
ancilla distance is at least the data-code distance under the assumptions of
the theorem.

\subsubsection{From row-parallel to addressable grid measurements}

The horizontal construction selects logical columns or column products, but
it repeats the same measurement in every logical row. Xu \emph{et al.}
obtain a selected grid pattern by adding a second HGP block $Q''$, called the
mask code. A vertical puncturing--augmenting homomorphism prepares selected
logical qubits of $Q'$ in $|+\rangle$, or prepares groups of them in logical
GHZ states, before the final horizontal-$Z$ measurement. Because the vertical
factor is represented by a transpose, the vertical homomorphism and its
physical CNOT reverse direction. Exchanging $X$ and $Z$ gives the dual
vertical-$X$ construction. Combining the two stages yields grid Pauli-product
measurements on products of disjoint row and column hyperedges.

\subsection{Requirements for an LP extension}

Equations~\eqref{eq:canonical-z}--\eqref{eq:canonical-x} suggest an LP analogue
of the HGP logical grid, but the extra fibre index $m$ creates additional
addressability requirements. A fault-tolerant construction should establish
all of the following rather than only an intended action on one family of
logical $Z$ operators:
\begin{enumerate}
  \item construct explicit $R$-linear ancillary complexes and chain maps, and
        identify their binary CNOT schedules after applying $\B$;
  \item compute both induced maps $\gamma_*$ and $\gamma^*$ on the entire
        logical space, including any residual sector outside the canonical
        fibres;
  \item select the desired $(i,j,m)$ logical qubits while proving that
        unmeasured canonical and residual logical information is preserved;
  \item bound the ancilla distance and the distance of the coupled or merged
        code, rather than inferring fault tolerance from the chain-map identity;
  \item give fault-tolerant ancilla preparation, repeated syndrome extraction,
        readout, decoding, and Pauli-frame procedures; and
  \item verify that $\B(\gamma_1)$ has bounded degree and admits a constant-depth
        or otherwise quantified physical CNOT schedule.
\end{enumerate}
The HGP construction supplies a concrete template---modify base complexes,
derive the chain map, calculate its complete logical action, and only then add
the preparation and readout protocol. The next step is to determine which of
these operations remain internal to the LP family over $R$ and which require a
separate mask or extractor construction.

\begin{thebibliography}{9}

\bibitem{ZhengEtAl2026}
H.~Zheng, G.~Zheng, L.~Jiang, and Q.~Xu,
``Logical Computation with Canonical Lifted Product Codes,''
\href{https://arxiv.org/abs/2607.28605v1}{arXiv:2607.28605v1} (2026).

\bibitem{XuEtAl2024}
Q.~Xu, H.~Zhou, G.~Zheng, D.~Bluvstein, J.~P.~Bonilla Ataides,
M.~D.~Lukin, and L.~Jiang,
``Fast and Parallelizable Logical Computation with Homological Product Codes,''
\href{https://arxiv.org/abs/2407.18490v1}{arXiv:2407.18490v1} (2024).

\end{thebibliography}

\end{document}
